
\subsubsection{2.1 Uncertainty Quantification for LLMs}

Semantic entropy \citep{kuhn2023semantic} computes uncertainty over clusters of semantically equivalent generations, achieving approximately 0.85 AUROC on TriviaQA. By using bidirectional entailment to group generations by meaning rather than surface form, semantic entropy captures linguistic invariances that token-level entropy does not address. However, this work---like most in the field---evaluates each benchmark independently without testing cross-benchmark transfer.

P(True) methods \citep{kadavath2022language} prompt models to predict the probability that their own outputs are correct. Calibration improves with model scale, but the approach requires explicit self-evaluation prompting and has not been evaluated across diverse benchmark families.

Token-level uncertainty measures including entropy and predictive variance have been explored for uncertainty estimation \citep{xiao2021hallucination}, but semantic entropy has shown improved performance by accounting for meaning-level rather than surface-level variation.

\subsubsection{2.2 Consistency-Based Detection}

SelfCheckGPT \citep{manakul2023selfcheckgpt} detects hallucinations by measuring consistency across multiple samples, achieving strong performance on WikiBio without requiring external knowledge. While methodologically distinct from entropy-based approaches, consistency methods face similar transfer questions. The clustering framework presented here could potentially extend to consistency-based methods.

\subsubsection{2.3 Benchmark Evaluation Paradigms}

Existing hallucination benchmarks---TriviaQA \citep{joshi2017triviaqa}, HaluEval \citep{li2023halueval}, FEVER \citep{thorne2018fever}---have been treated as interchangeable representatives of factual QA. However, no prior work has examined whether benchmarks cluster into distinct families based on the error processes they probe.

Domain adaptation literature \citep{bendavid2010theory} provides theoretical grounding: distribution shift degrades transfer, and distribution similarity predicts transferability. The present work operationalizes this theory for hallucination detection, using JS-divergence to measure uncertainty distribution similarity and hierarchical clustering to discover benchmark families.

